% ============================================================
% Capítulo: Materiales y Métodos (Bloque B.7)
% ============================================================
\chapter{Materiales y métodos}
\label{cap:materiales-metodos}

\section{Enfoque metodológico: Design Science Research}
\label{sec:mm-dsr}

SGB-SaaS se desarrolló siguiendo, de forma implícita durante el proyecto
y reconstruida aquí de forma explícita, la estructura de seis actividades
de \emph{Design Science Research} (DSR) de Peffers et al.
\citep{peffers2007dsrm} -- un proceso
iterativo pensado específicamente para investigación que produce un
artefacto (en este caso, software) como resultado, no solo conocimiento
descriptivo. La evaluación del artefacto se rige por las siete guías de
Hevner et al. \citep{hevner2004design} para la investigación en diseño
en sistemas de información, en particular la evaluación observacional
del artefacto desplegado en su entorno real de uso. Las seis actividades se mapean así contra la evolución real
y documentada del proyecto:

\begin{enumerate}[leftmargin=1.4cm]
  \item \textbf{Identificación del problema y motivación.} La fricción
    operativa real de la gestión bibliotecaria manual
    (\autoref{cap:introduccion}), delimitada en la fase de planificación
    del proyecto (Entrega~1A).
  \item \textbf{Definición de objetivos de una solución.} Las cuatro
    preguntas de investigación (RQ1--RQ4) y los cinco objetivos
    específicos (OE1--OE5) declarados en el \autoref{cap:introduccion}.
  \item \textbf{Diseño y desarrollo.} Construcción incremental del
    artefacto en sucesivas entregas -- Entrega~1B (primer módulo
    funcional), Entrega~3 (candidato estable con los ocho módulos y
    evidencia empírica inicial), Entrega Final (HEAD en \texttt{demo/interfaces-completas}, este
    documento) -- con las decisiones de arquitectura formalizadas como
    ADR y el modelo C4 (\autoref{cap:diseno-arquitectura}).
  \item \textbf{Demostración.} Ejecución de casos de uso reales de punta
    a punta contra el stack Docker completo, no simulados:
    \texttt{docs/\allowbreak mediciones/\allowbreak backend/\allowbreak 2026-07-29-flujo-prestamo-devolucion-multa-e2e.md}
    y
    \texttt{docs/\allowbreak mediciones/\allowbreak frontend/\allowbreak 2026-07-30-flujo-\allowbreak frontend-\allowbreak prestamos-\allowbreak reservaciones-\allowbreak multas-e2e.md}.
  \item \textbf{Evaluación.} El conjunto de instrumentos de medición
    descrito en la \autoref{sec:mm-instrumentos} de este mismo capítulo
    (k6, JaCoCo, Lighthouse, auditoría OWASP, y SUS pendiente de
    ejecución).
  \item \textbf{Comunicación.} Este mismo informe académico, el checklist
    de calidad de requisitos INCOSE~v4
    (\texttt{docs/checklists/incose2023-req.md}), la documentación
    arquitectónica versionada como código (\texttt{docs/arquitectura/}),
    y el artefacto de software archivado con DOI persistente en Zenodo
    (\autoref{cap:introduccion}).
\end{enumerate}

\section{Instrumentos de medición}
\label{sec:mm-instrumentos}

\begin{itemize}[leftmargin=1.2cm]
  \item \textbf{System Usability Scale (SUS)} \citep{brooke1996sus}: 10
    ítems tipo Likert de 5 puntos, alternando afirmaciones positivas y
    negativas, que producen una puntuación agregada de 0 a 100 por
    participante mediante la fórmula estándar del instrumento. Es el
    instrumento del Bloque~C.3 de esta guía; su estado de ejecución real
    -- todavía pendiente -- se declara con honestidad en la
    \autoref{sec:mm-poblacion}.
  \item \textbf{k6} (\texttt{k6/libros-listado-test.js},
    \texttt{k6/opts.js}): herramienta de prueba de carga HTTP, usada para
    medir la latencia del endpoint de catálogo bajo 50~VUs concurrentes
    en los escenarios \texttt{cache\_caliente}/\texttt{cache\_frio}. Su
    protocolo completo se describe en la \autoref{sec:mm-protocolo-perf}.
  \item \textbf{Lighthouse} (\texttt{docs/mediciones/lighthouse/lhci-20260731-0300.json}):
    auditoría automatizada de calidad del frontend, con puntuaciones
    reales obtenidas de 0,95 en Performance, 0,95 en Accessibility, 1,00
    en Best Practices y 0,82 en SEO sobre una escala de 0 a 1.
  \item \textbf{JaCoCo} (\texttt{jacoco-maven-plugin} en
    \texttt{backend-springboot/pom.xml}): cobertura de línea/rama/complejidad
    ciclomática sobre el backend, con objetivo declarado de la guía de
    $\geq$70\,\% de líneas para la Entrega Final. El resultado de cierre
    vigente es \textbf{56,01\,\%} de líneas (1613/2880) y
    \textbf{34,59\,\%} de ramas (247/714), por debajo del umbral de
    líneas -- tras un primer ciclo de trabajo dirigido específicamente a
    cinco clases de autenticación/autorización que partían de cobertura
    casi nula, una medición intermedia que confirmó que el porcentaje
    subió a 81,64\,\% tras el merge de los ocho módulos de dominio
    nuevos que casi triplicaron la base de código analizable
    (\texttt{docs/\allowbreak mediciones/\allowbreak jacoco/\allowbreak 2026-08-13-cobertura-jacoco-post-merge-8-modulos.md}),
    y una corrida de cierre posterior que registró un descenso a
    56,01\,\% tras la incorporación de funcionalidad adicional (ver
    \autoref{sec:res-cobertura} para el detalle completo).
  \item \textbf{Auditoría OWASP} (\texttt{scripts/owasp-audit.sh},
    \texttt{docs/mediciones/sec/}): verificación reproducible vía
    \texttt{curl} contra el stack Docker real de un subconjunto
    representativo de controles del OWASP Top~10:2021
    (\autoref{sec:mt-owasp}), con hallazgo y corrección documentados por
    control.
\end{itemize}

\section{Enfoque de métricas: Goal-Question-Metric (GQM)}
\label{sec:mm-gqm}

El bloque de evaluación de rendimiento de SGB-SaaS se estructuró,
retrospectivamente, según el patrón Goal-Question-Metric:

\begin{quote}
\textbf{Meta (\emph{Goal}).} Analizar el uso de caché Redis sobre el
endpoint de catálogo, con el propósito de evaluar su efecto sobre el
rendimiento observable (latencia), desde el punto de vista del equipo de
desarrollo, en el contexto de carga concurrente controlada (k6, 50~VUs).

\textbf{Preguntas (\emph{Questions}).}
\begin{itemize}[leftmargin=1cm]
  \item Q1. ¿Cuál es la latencia p95 del endpoint con caché caliente
    frente a caché frío?
  \item Q2. ¿La diferencia observada entre ambos escenarios es
    estadísticamente significativa, y de qué magnitud es el efecto?
  \item Q3. ¿Se cumplen los umbrales de latencia exigidos por la guía
    para cada escenario?
\end{itemize}

\textbf{Métricas (\emph{Metrics}).}
\begin{itemize}[leftmargin=1cm]
  \item M1. Percentiles p50/p90/p95/p99 de \texttt{http\_req\_duration}
    (ms) por escenario, sobre las 5 corridas agregadas
    (\texttt{scripts/perf-analysis.py}).
  \item M2. Estadístico y $p$-valor de la prueba de Wilcoxon de rangos
    con signo (pareada, sobre el p95 de cada corrida), y tamaño de
    efecto de Cliff's delta.
  \item M3. Cumplimiento binario (sí/no) de los umbrales exigidos: p95
    $<$200\,ms en caliente, p95 $<$500\,ms en frío.
\end{itemize}
\end{quote}

Los resultados reales obtenidos para M1--M3 están en
\texttt{docs/mediciones/perf/REPORT.md} y se transcriben en el
\autoref{cap:diseno-arquitectura}. Este mismo patrón GQM podría
extenderse a los otros bloques de evaluación (seguridad, cobertura,
usabilidad); no se desarrolla aquí para cada uno por disciplina de
alcance de este capítulo -- un GQM completo es suficiente para
documentar el enfoque de métricas adoptado, sin redundar en una
plantilla repetida para cada bloque.

\section{Población, muestreo y participantes}
\label{sec:mm-poblacion}

La población objetivo real de este proyecto es la comunidad de la
Universidad Técnica Estatal de Quevedo (UTEQ) en sus roles de dominio
(lector, bibliotecario) -- no una muestra sintética ni un panel externo
contratado. \citet{baltes2022sampling}, en su revisión crítica de
prácticas de muestreo en investigación de ingeniería de software,
advierten que muestras pequeñas y no aleatorizadas -- el perfil esperable
de cualquier estudio de este alcance, reclutado por conveniencia dentro
de una única institución -- limitan severamente la generalización de
cualquier conclusión más allá de la muestra concreta estudiada, y
recomiendan declarar explícitamente el método de reclutamiento en vez de
presentar el resultado como representativo de una población más amplia.

\textbf{Nota de honestidad -- estado real de ejecución.} La muestra SUS
\textbf{no se ha reclutado ni ejecutado todavía} al momento de escribir
este capítulo (\texttt{OBS-08} en
\texttt{docs/observaciones/OBSERVACIONES.md}: ``Pendiente -- en progreso,
asignada a Panama''). Lo que sí existe, ya definido y versionado antes de
la primera ejecución real (no de forma retroactiva), es el protocolo
completo: plantilla de consentimiento informado
(\texttt{docs/etica/consentimientos/plantilla.md}), técnica de muestreo
por conveniencia dentro de la comunidad UTEQ, tamaño mínimo de muestra
$N\geq15$ participantes exigido por la guía, y anonimización mediante
código de participante (\texttt{P01}, \texttt{P02}, \ldots, ver
\texttt{docs/etica/ETHICS.md}, sección~iii) -- nunca nombre ni correo. No
se declara una muestra ya obtenida para no incurrir en el mismo tipo de
afirmación no verificada que este documento evita de forma sistemática en
el resto de sus capítulos.

\section{Protocolo experimental de rendimiento (replicable)}
\label{sec:mm-protocolo-perf}

El bloque de rendimiento sigue un protocolo completamente reproducible,
documentado con el comando exacto ejecutado:

\begin{lstlisting}[language=bash]
docker run --rm --network sgb-saas_default \
  -v "$(pwd)/k6:/scripts" -v "$(pwd)/docs/mediciones/perf:/out" \
  grafana/k6 run --out json=/out/k6-run{N}.json /scripts/libros-listado-test.js
\end{lstlisting}

Se ejecutó 5 veces de forma independiente (el mínimo exigido por la
guía), cada corrida cubriendo ambos escenarios de forma secuencial
(\texttt{cache\_caliente}: siempre \texttt{page=0}, para forzar
repetidos aciertos de caché sobre la misma clave; \texttt{cache\_frio}:
página elegida por un PRNG determinista (\texttt{mulberry32}, semilla
fija~42) en cada petición, para forzar fallo de caché en cada una), con
un perfil de carga común de 10s de \emph{ramp-up} a 50~VUs, 30s
sostenidos, 10s de \emph{ramp-down}
(\texttt{k6/opts.js}). Autenticación real vía \texttt{POST /api/auth/login}
en el \texttt{setup()} del script, reutilizando el mismo
\texttt{accessToken} entre VUs. Los cinco archivos de datos crudos
(\texttt{k6-run1.json}--\texttt{k6-run5.json}, formato NDJSON de k6) están
versionados en \texttt{docs/mediciones/perf/}, permitiendo reanalizar los
mismos datos sin repetir la carga. El detalle completo -- entorno exacto
(versiones de Docker, Java, PostgreSQL, Redis, k6), resultados por
corrida y agregados -- está en \texttt{docs/mediciones/perf/REPORT.md} y
se resume en el \autoref{cap:diseno-arquitectura}.

\section{Análisis estadístico}
\label{sec:mm-estadistica}

El análisis agregado de las 5 corridas (\texttt{scripts/perf-analysis.py})
calcula, por escenario, sobre la métrica \texttt{http\_req\_duration}: la
media, la mediana, la desviación típica ($\sigma$), el intervalo de
confianza al 95\,\% de la media, los percentiles p50/p90/p95/p99, la tasa
de error HTTP $\geq$500 y el \emph{throughput} (peticiones/segundo). Sobre
la comparación pareada \texttt{cache\_caliente} vs.\ \texttt{cache\_frio}
(p95 de cada una de las 5 corridas, no el \emph{pool} de peticiones
agregado, precisamente porque las corridas son las mismas 5 sesiones de
carga y no muestras independientes) se aplica la prueba de Wilcoxon de
rangos con signo (método exacto, vía \texttt{scipy.stats.wilcoxon}) y el
tamaño de efecto de Cliff's delta. Los resultados reales obtenidos --
Wilcoxon $p=0{,}0625$, Cliff's delta $=-0{,}68$ -- y su interpretación
correcta en función de la potencia estadística del tamaño de muestra
mínimo se desarrollan en el \autoref{cap:amenazas-validez} (validez de
conclusión).

\section{Determinismo y semillas}
\label{sec:mm-semillas}

Para que los resultados sean reproducibles bit a bit donde el diseño
experimental usa aleatoriedad, se fijaron semillas explícitas en los dos
únicos puntos del repositorio donde se generan números pseudoaleatorios
con fines de medición (verificado por búsqueda directa en el código):

\begin{itemize}[leftmargin=1.2cm]
  \item \textbf{Selección de página en el escenario \texttt{cache\_frio}}
    (\texttt{k6/libros-listado-test.js}): PRNG \texttt{mulberry32} con
    semilla fija \textbf{42}, de modo que la secuencia de páginas pedidas
    en cada corrida es la misma si se re-ejecuta el script.
  \item \textbf{Intervalo de confianza al 95\,\% por \emph{bootstrap} del
    percentil p95} en el gráfico comparativo
    (\texttt{scripts/perf-analysis.py},
    \texttt{BOOTSTRAP\_SEED = 42}): 2000 réplicas de remuestreo con
    reemplazo, con la misma semilla \textbf{42} -- un percentil, a
    diferencia de la media, no tiene una fórmula cerrada de intervalo de
    confianza, de ahí el uso de \emph{bootstrap}.
\end{itemize}

Ninguna otra parte del repositorio (scripts de prueba, seed de base de
datos, generación de datos de ejemplo) usa aleatoriedad con fines de
medición que requiera declarar una semilla adicional.
